\section{Seguridad, Pruebas de Rendimiento y Despliegue con Docker}
\subsection{Auditoría de seguridad OWASP en SGROAS}

Se realizó una auditoría de seguridad sobre la aplicación MVC usando el OWASP Top 10:2021 (A01--A07) más XSS. La
Tabla~\ref{tab:owasp} resume vulnerabilidad, vector de ataque, contramedida implementada y evidencia en el
repositorio (directorio \texttt{docs/mediciones/sec/}).

\begin{table}[H]
\centering
\footnotesize
\setlength{\tabcolsep}{4pt}
\begin{tabular}{@{}p{2.4cm}p{2.9cm}p{3.3cm}p{2.6cm}@{}}
\toprule
\textbf{Vulnerabilidad} & \textbf{Vector de ataque} & \textbf{Contramedida implementada} & \textbf{Evidencia} \\
\midrule
A01 Control de acceso & Acceso no autorizado a /api/conductores & Reglas por rol en SecurityConfig (ROLE\_ADMIN, ROLE\_COORDINADOR) & A01-acceso.txt \\
A02 Fallo criptográfico & Robo de credenciales/token & BCrypt, JWT HS256, cookie HttpOnly+Secure+SameSite & A02-criptografia.txt \\
A03 Inyección & SQL/JPQL malicioso en login & Validación @Valid y consultas JPA parametrizadas & A03-inyeccion.txt (HTTP 422) \\
A05 Malconfiguración & Cabeceras inseguras & HSTS, X-Frame-Options, CSP, nosniff & A05-headers.txt \\
A07 Fallos de autenticación & Fuerza bruta en login & Rate limiter (6 intentos/60 s) → HTTP 429 & A07-rate-limit.txt \\
A09 Monitoreo & Falta de logs & Logging estructurado del filtro JWT & A09-logs.txt \\
XSS & Script malicioso en el navegador & CSP (default-src 'self') y sanitización de Angular & SecurityConfig + frontend \\
\bottomrule
\end{tabular}
\caption{Resumen de la auditoría OWASP aplicada al PFC.}
\label{tab:owasp}
\end{table}

\textbf{A01 — Control de acceso roto.} El vector consiste en invocar \texttt{GET /api/conductores} sin token o con
un rol insuficiente. La contramedida es la configuración declarativa de seguridad: solo
\texttt{ROLE\_ADMIN} y \texttt{ROLE\_COORDINADOR} pueden operar sobre los conductores; la evidencia
\texttt{A01-acceso.txt} registra la respuesta 401/403 obtenida.

\textbf{A02 — Fallos criptográficos.} Los vectores son credenciales débiles, tokens en tránsito o en almacenamiento
inseguro. Las contramedidas son: cifrado BCrypt para contraseñas, firma HS256 del JWT con clave de mínimo 32
caracteres, y cookie \texttt{HttpOnly, Secure, SameSite=Strict}. El payload firmado contiene 7 claims
(RFC 7519), como consta en \texttt{A02-criptografia.txt}.

\textbf{A03 — Inyección.} El vector típico es un payload \texttt{' OR 1=1 --} en el login. La contramedida es doble:
la validación Jakarta (\texttt{@Email}, \texttt{@NotBlank}) rechaza el payload antes de llegar a la BD con HTTP 422
(ProblemDetails), y las consultas se construyen con JPA/HQL parametrizado, sin concatenar entrada de usuario, tal
como se verificó en \texttt{A03-inyeccion.txt}.

\textbf{A05 — Mala configuración de seguridad.} El vector son las cabeceras HTTP inseguras. SGROAS implementa
\texttt{Strict-Transport-Security} (1 año, includeSubDomains), \texttt{X-Frame-Options: DENY},
\texttt{X-Content-Type-Options: nosniff}, \texttt{Content-Security-Policy: default-src 'self'} y
\texttt{Cache-Control: no-cache}, verificados en \texttt{A05-headers.txt}.

\textbf{A07 — Fallos de autenticación.} El vector es la fuerza bruta sobre \texttt{/api/auth/login}. La
contramedida es \texttt{LoginRateLimiter}, un limitador en memoria (ConcurrentHashMap) que permite 6 intentos
fallidos por ventana de 60 segundos y responde HTTP 429. La evidencia \texttt{A07-rate-limit.txt} muestra la
secuencia de intentos y el código 429 obtenido.

\textbf{XSS — Cross-Site Scripting.} El vector es la ejecución de scripts en el navegador mediante datos
maliciosos. La contramedida principal es la \texttt{Content-Security-Policy} con \texttt{default-src 'self'} y
\texttt{script-src 'self'}, además del saneamiento automático que realiza Angular por defecto al interpolar
valores en plantillas.

\subsection{Prueba de carga con k6 (Apache Bench)}

La metodología de prueba de carga del PFC se realizó con \textbf{k6} (también se documentó el uso de Apache Bench
con \texttt{ab -n <total> -c <concurrentes> http://localhost:8080/api/conductores}). La config fue: \textbf{50
usuarios virtuales concurrentes} durante 30 segundos, con \textit{ramp-up}, contra el endpoint representativo
\texttt{GET /api/conductores} (protegido y con caché de Redis caliente). Se ejecutaron 3 corridas independientes
(archivos \texttt{k01}, \texttt{k02}, \texttt{k03}). Los resultados se resumen en la Tabla~\ref{tab:perf}.

\begin{table}[H]
\centering
\small
\begin{tabular}{@{}lrrrrr@{}}
\toprule
\textbf{Métrica} & \textbf{Corrida 1} & \textbf{Corrida 2} & \textbf{Corrida 3} & \textbf{Media} \\
\midrule
Peticiones & 1466 & 1501 & 1501 & 1489 \\
p50 (mediana, ms) & 17,88 & 10,48 & 9,13 & --- \\
p95 (ms) & 173,15 & 34,82 & 36,32 & 81,43 \\
p99 (ms) & 251,54 & 155,99 & 156,97 & 188,16 \\
Throughput (req/s) & --- & --- & --- & 49,64 \\
Errores & 0 & 0 & 0 & 0 \% \\
\bottomrule
\end{tabular}
\caption{Resultados de la prueba de carga (k6, 50 VUs por 30 s, commit 62bf8fa).}
\label{tab:perf}
\end{table}

La media de tiempo de respuesta fue de \textbf{22,93 ms}, con un \textbf{p95 de 81,43 ms}, muy por debajo del
umbral objetivo declarado en \texttt{k6/opts.js} (p95 < 200 ms). El \textbf{throughput alcanzó 49,64 req/s} y la
\textbf{tasa de error fue 0,00 \%}. La corrida 1 es la más alta (p95 de 173,15 ms) por el arranque de la caché
caliente; las corridas 2 y 3 estabilizan el p95 por debajo de 40 ms. \textbf{Criterio: Cumplido.}

\subsection{Containerización con Docker Compose}

Docker Compose permite desplegar la aplicación con entornos reproducibles y aislamiento de dependencias. El
\texttt{docker-compose.yml} del PFC define \textbf{tres servicios} en una red bridge \texttt{sgroas-network}:
\begin{itemize}
    \item \textbf{postgres} (postgres:18, puerto 5433:5432, volumen \texttt{sgroas\_postgres\_data}).
    \item \textbf{redis} (redis:7, puerto 6379) para caché (TTL y blacklist de JWT).
    \item \textbf{backend} (build local, puerto 8080:8080, con variables de entorno para BD, Redis y JWT).
\end{itemize}

El \texttt{Dockerfile} es multi-etapa: fase de build con \texttt{eclipse-temurin:21-jdk} (\texttt{./mvnw clean
package -DskipTests}) y etapa final con \texttt{eclipse-temurin:21-jre} que ejecuta \texttt{java -jar app.jar}. Las
variables de entorno (\texttt{SPRING\_DATASOURCE\_URL}, \texttt{SPRING\_DATA\_REDIS\_HOST},
\texttt{APP\_JWT\_SECRET}, \texttt{APP\_JWT\_EXPIRATION\_MS}, etc.) se inyectan desde el entorno y \textbf{nunca
deben estar en el código fuente}: en el repositorio solo existe \texttt{.env.example} como plantilla, y la
configuración de desarrollo usa valores por defecto que en producción se sobrescriben. Un despliegue real usaría
secretos con HashiCorp Vault o los secretos del orquestador (Docker secrets / GitHub Actions secrets). El
despliegue se verifica con \texttt{docker compose up -d} y \texttt{docker compose ps} (todos los servicios en
estado \texttt{Up}).

\subsection{Reflexión sobre el proceso de aprendizaje}

\textbf{Tejada Bajaña Luis Alejandro.} A lo largo del desarrollo del PFC, la lección técnica más importante fue
comprender que la seguridad no es una capa opcional sino una restricción transversal que condiciona el diseño: la
elección de JWT stateless con refresh token en Redis, el rate limiting de 6 intentos/60 s en el login y las
contramedidas frente a A01 y A03 debieron decidirse al inicio, no al final. El error más costoso fue asumir que la
validación en el frontend bastaba; la auditoría demostró que el backend debía validar todo payload con Jakarta
Validation y devolver ProblemDetails 422 antes de tocar la base de datos. En un proyecto futuro validaría y
auditaría la seguridad de forma continua (security-by-design) y aplicaría los principios de \textit{Release It!}
(timeouts, reintentos y fallas parciales) al consumo de servicios externos desde la primera iteración.

\textbf{Alava Alvarado Jean Pierre.} La experiencia de la evaluación cruzada me permitió entender el valor del
código documentado y de las decisiones registradas en ADR: al revisar un proyecto ajeno, los ADR y la colección
Postman fueron la mejor puerta de entrada para comprender por qué se eligió el monolito con Spring Boot y cómo se
resolvió la revocación de JWT con la blacklist de Redis. El error que identifiqué como más costoso en términos de
comunicación es no mantener actualizada la versión del frontend en la documentación (el README indica Angular 17
mientras el \texttt{package.json} declara Angular 20), lo que genera fricción en el equipo. En el futuro priorizaría
registrar la trazabilidad éxito-criterio desde el inicio (como la matriz \texttt{docs/trazabilidad/matriz.csv}) y
ejecutar la prueba de carga con caché fría y caliente desde los primeros entregables, para detectar pronto los
puntos de degradación.